\section{De Onderneming als Organisatie}
De volgende hoofdstukken gaan over hoeveel de producenten gaan produceren en tegen welke prijs ze die productie gaan verkopen. In dit hoofdstuk: aan welke voorwaarden moet voldaan zijn om een maximale winst te hebben.

\subsection{Waarom Ondernemingen?}
eenmanszaak, BVBA, NV, blablabla...

Waarom zijn er nu ondernemingen en gaat niet alles via individuele spelers die hun output op de markt verkopen? Een individu neemt dan een taak op zich die in ons systeem een stap in het proces van een onderneming is.

Dit komt door de transactiekosten dat die werkwijze met zich meebrengt. Op termijn zullen er contracten tussen individuen ontstaan, maar deze contracten kunnen onmogelijk op alle eventualiteiten inspelen $\rightarrow$ mensen gaan zich groeperen, autoriteit aan \'e\'en individu afstaan om deze contracten dan te internaliseren binnen het bedrijf dat ontstaat.

\subsection{Winstmaximalisatie}
Ondernemingen zijn voornamelijk ge\"interesseerd in het maximaliseren van hun winst. Deze wordt beschreven met volgende functie:
\begin{equation}
	W(q) = TO(q) - TK(q)
\end{equation}

Deze bestaat dus uit de opbrengsten min de kosten. Deze kosten zijn de economische kosten, dus ook zaken zoals opportuniteitskosten worden in rekening gebracht. Aangezien $W$ in functie staat van $q$, zullen we de $q$ zoeken die $W$ maximaliseert. 

\subsubsection{Ontvangstenfuncties}
Volgende functies zijn relevant aan de opbrengstenzijde:
\begin{subequations}
\begin{align}
	TO(q) &= q \times q(p)\\
    GO(q) &= \frac{TO(q)}{q} = \frac{q \times p(q)}{q} = p(q)\\
    MO(q) &= \frac{\Delta TO(q)}{\Delta q} \rightarrow MO(q) = \frac{\mathrm{d} TO(q)}{\mathrm{d}q}
\end{align}
\end{subequations}
$p(q)$ is daarbij de inverse vraagfunctie. Merk op dat wanneer $q$ veranderd, $TO$ twee effecten zal ondergaan(figuur~\ref{fig:GOenTO}): een effect van de veranderde hoeveelheid en een verandering in de andere richting van de prijs. Het totale effect zal van de $\varepsilon_{V}^p$ afhangen. Als je naar het middelpunt beweegt zal $TO$ stijgen, als je weg beweegt dalen. Dit vertaalt zich ook in het snijden van $MO$ met de x-as in het punt waar die $q$ ook bij het punt hoort waar $|\varepsilon_{V}^p| = 1$.

\begin{figure}[htbp]
	\centering
	\includegraphics[scale=0.4]{Images/white.png}
	\caption{Inkomens- en substitutie-effect grafisch weergegeven}
	\label{fig:GOenTO}
\end{figure}

$GO$ en $MO$ kunnen we ook afleiden als we $TO$ als gegeven hebben. In elke $q$ is de waarde van $MO$ gelijk aan de afgeleide van $TO$ in dat punt. In elke $q$ is $GO$ gelijk aan $TO$ in dat punt gedeeld door die $q$.


\subsubsection{Kostenfuncties}
=verband tussen produciehoeveelheid en de kosten om die hoeveelheid te kunnen produceren. Analoog aan de opbrengstenfuncties hebben we nu $TK$, $GK$ en $MK$, waarbij $MK$ altijd positief moet zijn aangezien er altijd meer geld nodig zal zijn voor extra inputs. We gaan er voorlopig van uit dat $TK$ niet lineair is. Waarom juist zien we in hoofdstuk~\ref{sec:productieEnKosten}.

Het gevolg van die niet-lineariteit is dat zowel $MK$ en $GK$ eerst dalen en daarna terug stijgen. $MK$ snijdt $GK$ ook in zijn laagste punt en ligt voor dat punt onder $GK$. Ook hier kunnen $MK$ en $GK$ afgeleid worden uit de curve van $TO$ op een analoge manier als voorheen.


\subsubsection{Optimale Productie}
Er zijn twee voorwaarden om in een punt van optimale productie te zijn:
\begin{enumerate}
	\item Zoek de $q^*$ met de hoogste winst, waarbij $q^* > 0$
    \item \textbf{sluitingsregel}: ga na of de winst in deze $q^*$ ook effectief positief is
\end{enumerate}

Het zoeken naar een optimale $q^*$ kunnen we op twee manieren doen: kijken waar $TW$ het grootst is, met $ TW(q) = TO(q) - TK(q)$. Een andere, meer gebruikte manier is zoeken naar het punt waar $MW(q) = 0$ met $MW(q) = MO(q) - MK(q)$. Bij deze laatste methode moet je wel zorgen dat:
\begin{equation}
	\left \{\begin{array}{l l}
    MO(q) > MK(q) \text{ voor } q < q^*\\
    MO(q) < MK(q) \text{ voor } q > q^*
  \end{array} \right.
\end{equation}

We weten nu dat in $q^*$ de winst maximaal is voor alle positieve productieniveaus, maar het kan nog altijd zijn dat je daar verlies maakt ($q^*$ is de prod. met het minimale verlies).

$\rightarrow$ controleren of $TW(q)$ in $q^*$ positief is (=\textbf{sluitingsregel})



